\begingroup
\tcbset{myboxstyle/.append style={breakable=false}}
\Needspace{18\baselineskip}
\subsection{Основные абсолютно непрерывные распределения}

\subsubsection{Равномерное распределение}

\begin{definition}
    Случайная величина $\xi$ имеет \textbf{равномерное распределение}
    на отрезке $[a;b]$, где $a<b$, если её плотность имеет вид
    \[
        p_\xi(t)=\frac{1}{b-a}\,I_{[a;b]}(t)
        =\begin{cases}
            \dfrac{1}{b-a},&a\leq t\leq b,\\
            0,&t\notin[a;b].
        \end{cases}
    \]
    Обозначение: $\xi\sim\Uniform[a;b]$.
    Здесь $I_A(t)=1$ при $t\in A$ и $I_A(t)=0$ при $t\notin A$.
\end{definition}

\begin{center}
    \begin{tikzpicture}[x=1.4cm,y=1.8cm,>=Stealth]
        \draw[->] (-0.5,0)--(4.6,0) node[right] {$t$};
        \draw[->] (0,-0.2)--(0,1.6) node[above] {$p_\xi(t)$};
        \draw[dashed] (0,1)--(1,1);
        \draw[dashed] (1,0)--(1,1) (3,0)--(3,1);
        \draw[blue,very thick] (-0.45,0)--(1,0) (1,1)--(3,1) (3,0)--(4.4,0);
        \filldraw[blue,fill=white] (1,0) circle (1.5pt) (3,0) circle (1.5pt);
        \fill[blue] (1,1) circle (1.5pt) (3,1) circle (1.5pt);
        \node[below left] at (0,0) {$0$};
        \node[below] at (1,0) {$a$};
        \node[below] at (3,0) {$b$};
        \node[left] at (0,1) {$\frac1{b-a}$};
    \end{tikzpicture}
\end{center}

\newpage
\subsubsection{Экспоненциальное распределение}

\begin{definition}
    Случайная величина $\xi$ имеет \textbf{экспоненциальное распределение}
    с параметром $\lambda>0$, если её плотность имеет вид
    \[
        p_\xi(t)=\lambda e^{-\lambda t}I_{[0;+\infty)}(t)
        =\begin{cases}
            \lambda e^{-\lambda t},&t\geq0,\\
            0,&t<0.
        \end{cases}
    \]
    Обозначение: $\xi\sim\operatorname{Exp}(\lambda)$.
\end{definition}

\[
    \int_{-\infty}^{+\infty}p_\xi(t)\,\diff t
    =\int_0^{+\infty}\lambda e^{-\lambda t}\,\diff t
    =\left[-e^{-\lambda t}\right]_0^{+\infty}=1.
\]

\begin{center}
    \begin{tikzpicture}[x=1.4cm,y=1.8cm,>=Stealth]
        \draw[->] (-0.5,0)--(4.6,0) node[right] {$t$};
        \draw[->] (0,-0.2)--(0,1.6) node[above] {$p_\xi(t)$};
        \draw[blue,very thick] (-0.45,0)--(0,0);
        \draw[blue,very thick,domain=0:4.3,samples=100]
            plot (\x,{exp(-\x)});
        \filldraw[blue,fill=white] (0,0) circle (1.5pt);
        \fill[blue] (0,1) circle (1.5pt);
        \node[left] at (0,1) {$\lambda$};
        \node[below left] at (0,0) {$0$};
    \end{tikzpicture}
\end{center}

\subsubsection{Нормальное распределение}

\begin{definition}
    Случайная величина $\xi$ имеет \textbf{нормальное распределение}
    с параметрами $a\in\R$ и $\sigma^2>0$, если её плотность равна
    \[
        p_\xi(t)=\frac{1}{\sqrt{2\pi}\,\sigma}
        \exp\left(-\frac{(t-a)^2}{2\sigma^2}\right),
        \qquad t\in\R,\quad\sigma>0.
    \]
    Обозначение: $\xi\sim N(a;\sigma^2)$.
    Распределение $N(0;1)$ называется \textbf{стандартным нормальным};
    его плотность равна
    \[
        p_\xi(t)=\frac{1}{\sqrt{2\pi}}e^{-t^2/2}.
    \]
\end{definition}

\newpage
\begin{lemma}[Интеграл Гаусса]
    \[
        \int_{-\infty}^{+\infty}e^{-x^2}\,\diff x=\sqrt\pi.
    \]
\end{lemma}

\begin{Proof}
    Обозначим интеграл через $C$. Он положителен и конечен:
    при $|x|\geq1$ имеем $e^{-x^2}\leq e^{-|x|}$.
    Тогда
    \[
        \begin{aligned}
            C^2
            &=\left(\int_{-\infty}^{+\infty}e^{-x^2}\,\diff x\right)
              \left(\int_{-\infty}^{+\infty}e^{-y^2}\,\diff y\right)\\
            &=\int_{-\infty}^{+\infty}\int_{-\infty}^{+\infty}
              e^{-(x^2+y^2)}\,\diff x\,\diff y
              =\iint_{\R^2}e^{-(x^2+y^2)}\,\diff x\,\diff y.
        \end{aligned}
    \]
    Перейдём к полярным координатам:
    \[
        x=r\cos\varphi,\qquad y=r\sin\varphi,
        \qquad r\geq0,\quad0\leq\varphi<2\pi.
    \]
    Якобиан замены равен
    \[
        \left|\frac{\partial(x,y)}{\partial(r,\varphi)}\right|
        =\left|\begin{vmatrix}
            \cos\varphi&-r\sin\varphi\\
            \sin\varphi&r\cos\varphi
        \end{vmatrix}\right|=r.
    \]
    Поэтому
    \[
        \begin{aligned}
            C^2&=\int_0^{+\infty}\int_0^{2\pi}
                  e^{-r^2}r\,\diff\varphi\,\diff r\\
               &=2\pi\int_0^{+\infty}r e^{-r^2}\,\diff r
                 =\pi\int_0^{+\infty}e^{-u}\,\diff u=\pi,
                 \qquad u=r^2.
        \end{aligned}
    \]
    Так как $C>0$, получаем $C=\sqrt\pi$.
\end{Proof}

\begin{Proof}[Нормировка нормальной плотности]
    При замене $y=(t-a)/(\sqrt2\,\sigma)$ имеем
    $\diff t=\sqrt2\,\sigma\,\diff y$. Следовательно,
    \[
        \begin{aligned}
            \int_{-\infty}^{+\infty}p_\xi(t)\,\diff t
            &=\frac{1}{\sqrt{2\pi}\,\sigma}
              \int_{-\infty}^{+\infty}
              e^{-(t-a)^2/(2\sigma^2)}\,\diff t\\
            &=\frac1{\sqrt\pi}\int_{-\infty}^{+\infty}
              e^{-y^2}\,\diff y=1.
        \end{aligned}
    \]
\end{Proof}
\endgroup
